%======================================================================
\section{Ruins and Retrospective Openness}\label{sec:ruins}
%======================================================================

A completed surface hides many of the decisions required to produce it. Plaster conceals the lath and the wiring, a finished interface conceals the code paths behind it, a consolidated statute conceals the amendments and compromises from which it was assembled. Decay reverses the concealment. A ruin exposes layers, repairs, joints, and incompatible historical periods; it shows where one campaign of building stopped and another began, where a wall was patched with salvaged material, where an opening was blocked. The ruin is thus an artifact whose internal dependency structure becomes more observable as its functional integrity declines.

\emph{Instantiation.} Here $\Omega$ is the physical and informational state of an aging artifact; $A$ is the set of redirections still possible; $\Sigma$ is what can be observed of its construction; $\mathcal{K}$ is the cost of intervention; $\mathcal{W}$ is the authority to repair, alter, or conserve; and $\sim_{\mathcal{R}}$ is the identity that conservation seeks to maintain.

\subsection{Observability and usability}

Let $G=(V,\mathcal{E})$ be the artifact's dependency graph, and for each edge $e\in\mathcal{E}$ let $h_e(t)\in[0,1]$ be the degree to which $e$ is concealed by covering layers at time $t$, and $\mathrm{doc}_e\in[0,1]$ the degree to which $e$ is independently recorded. Define the observability of the artifact's construction as
\begin{equation}
\Theta(t)\ =\ \frac{1}{|\mathcal{E}|}\sum_{e\in\mathcal{E}}\max\big\{1-h_e(t),\ \mathrm{doc}_e\big\},
\end{equation}
and let its usability $\Psi(t)$ be an increasing function of the integrity of its layers, $\Psi(t)=\psi(\bar h(t))$, where $\bar h$ is the mean covering integrity and $\psi$ is strictly increasing with $\psi'>0$. The assumption that usability rises with covering integrity expresses the fact that the layers which conceal are often the same layers which protect, enclose, and finish.

\begin{proposition}[The paradox of the ruin]\label{prop:ruin}
If the artifact is undocumented ($\mathrm{doc}_e=0$ for all $e$) and its mean covering integrity strictly declines ($\dot{\bar h}<0$), then
\[
\frac{d\Theta}{dt}>0\qquad\text{while}\qquad\frac{d\Psi}{dt}<0 .
\]
If instead $\mathrm{doc}_e\geq d_0$ for all $e$, then $\Theta(t)\geq d_0$ at every $t$, independently of $\Psi$.
\end{proposition}

\begin{proof}
With $\mathrm{doc}_e=0$, $\Theta(t)=1-\bar h(t)$, so $\dot\Theta=-\dot{\bar h}>0$ and $\dot\Psi=\psi'(\bar h)\dot{\bar h}<0$. With documentation, each summand is at least $d_0$.
\end{proof}

In the undocumented case the artifact moves along a decreasing frontier in the $(\Theta,\Psi)$ plane: it becomes easier to interpret as constructed precisely when it becomes harder to redirect. Understanding arrives when the capacity to act on it is departing, which is the temporal structure of \eqref{eq:defer-decay} in which information outruns the survival of options. Riegl identified the aesthetic side of this paradox. He distinguished the \emph{age value} of a monument, the immediate, widely accessible appreciation of the visible traces of time and decay, from its \emph{historical value} as a document of a particular moment and from the \emph{use value} that demands it be maintained in working order, and he observed that these values conflict: age value counsels letting decay proceed, while use value counsels repair that erases the traces age value prizes \cite{riegl1982}. Benjamin drew the analogy between ruins in the realm of things and allegory in the realm of thought, and read the fragment as disclosing what the finished whole suppresses: the contingency of its assembly and its exposure to history \cite{benjamin1998}. In both accounts the ruin reveals something the completed object had hidden, and in both the revelation comes at the object's expense. In terms of the scoped definition of completion, the ruin is open within an epistemic scope, its construction history, and closed within an operational scope, its original function. Decay moves these two scopes in opposite directions.

\subsection{The prospective ruin}

Proposition~\ref{prop:ruin} also indicates how the paradox can be escaped. Documentation decouples observability from decay. A well-designed system should acquire some epistemic properties of a ruin before suffering ruinous damage. It should expose its assembly, so that its dependency structure is legible while it still functions. It should preserve rejected alternatives, recording the branches that closure removed, so that a later continuer can reconsider them without reconstructing them. It should record repairs, so that the history of its revisions is part of its present state rather than something recovered by archaeology. And it should reveal intervention points, marking where it can be entered without destruction. Call such an artifact a \emph{prospective ruin}: one that already displays, in working order, the information that decay would otherwise disclose too late.

The prospective ruin is more demanding than transparency, provenance, or ordinary documentation, and the difference is worth stating. Transparency discloses the current structure; provenance records where components came from; documentation describes how the artifact works now. A ruin discloses more: superseded structures still embedded in the present one, repairs that record past failures, and blocked openings that mark continuations once available and later withdrawn. A prospective ruin therefore preserves not only the present dependency graph but its history and its counterfactuals, including the branches that were considered and rejected, with the reasons for rejecting them. Design-rationale methods that record the questions a design faced, the options considered, and the criteria by which they were chosen are the closest existing practice \cite{maclean1991}. They are rarely maintained, and Observation~\ref{obs:measurability} predicts why: rejected options leave no countable trace.

These practices correspond to the continuation infrastructure of Section~\ref{sec:inert} and to the open case of Section~\ref{sec:institutions}. A design rationale that records rejected options is a ledger of possibilities in the sense of Section~\ref{sec:option}; a version history is a record of repairs; an architectural drawing set kept with the building, a service panel labeled by circuit, a code base with its tests and decision records, a statute published with its legislative history, are all ways of holding $\mathrm{doc}_e$ high so that $\Theta$ need not wait for $\Psi$ to fall. The ruin shows retrospectively what latent incompletion requires prospectively: that the artifact be interpretable as made, and therefore as remakeable, while it can still be remade.

%======================================================================
\section{The Ethics of Reopening}\label{sec:reopening}
%======================================================================

An objection is implicit in any celebration of continuability. If every settlement remains reopenable, then promises, rights, and protections become unreliable. People plan their lives around closures: a pension formula, a property boundary, a legal protection, a court's final judgment. A system in which such closures can be revisited at will offers no ground to stand on, and the ground it removes is most valuable to those with the fewest alternatives. The objection is correct against a naive version of the thesis, and answering it requires two distinctions.

\subsection{Revisability and revocability}

\begin{definition}
A closure on component $j$ is \emph{revisable} if it can be altered through a specified procedure that requires evidence of strength at least a threshold $\theta_j$ and is available on stated terms to stated parties. It is \emph{revocable} if some authority can withdraw it unilaterally, without meeting any evidential threshold.
\end{definition}

A protected commitment can therefore remain institutionally revisable without becoming arbitrarily defeasible. A constitutional right that can be amended only by supermajorities after extended deliberation is revisable and not revocable. A benefit that an agency can withdraw by internal memorandum is revocable. Continuability as defended here concerns revisability. Revocability is the privilege of the authority whose concentrated rows in $Q$ Proposition~\ref{prop:revision-capacity} warned against.

The threshold $\theta_j$ should depend on four quantities: the severity of error if the closure is mistaken, the quality of evidence that could bear on it, the reliance interests the existing arrangement has created, and the distribution of harms that reopening would produce. Write $\theta_j=\theta(\text{severity}_j,\text{evidence}_j,\text{reliance}_j,\text{distribution}_j)$, decreasing in severity (grave errors should be easier to correct) and increasing in reliance (arrangements others depend upon should be harder to disturb). The difficult cases are those in which these pull in opposite directions, and there the threshold function is where the normative argument must be made explicit rather than left to whoever controls the procedure.

\subsection{Symmetric reopening rules and repeated challenge}

The second distinction concerns symmetry. It may seem fair that any party may seek to reopen any settlement under the same threshold. The following result shows why such symmetry reproduces inequality.

Model a challenge to a correct settlement as an independent draw of evidence strength $\epsilon\sim N(0,1)$, the noise that any finite inquiry produces, with reopening if $\epsilon\geq\theta$. The per-challenge probability of wrongly reopening is $p(\theta)=1-\Phi(\theta)$.

\begin{proposition}[Attrition of settlements: independent-challenge benchmark]\label{prop:attrition}
An agent able to mount $n$ independent challenges to a correct settlement under a fixed threshold $\theta$ overturns it with probability $1-\Phi(\theta)^n$, which tends to one as $n\to\infty$. To hold the probability of wrongful overturn below $\alpha\in(0,1)$ against $n$ challenges, the threshold must satisfy
\begin{equation}\label{eq:threshold-n}
\theta\ \geq\ \Phi^{-1}\!\big((1-\alpha)^{1/n}\big),
\end{equation}
which grows without bound in $n$, asymptotically as $\sqrt{2\log n}$.
\end{proposition}

\begin{proof}
The settlement survives $n$ independent challenges with probability $\Phi(\theta)^n$, which tends to zero for fixed $\theta$. Requiring $1-\Phi(\theta)^n\leq\alpha$ gives $\Phi(\theta)\geq(1-\alpha)^{1/n}$, which is \eqref{eq:threshold-n}. Since $(1-\alpha)^{1/n}=1-\frac{-\log(1-\alpha)}{n}+O(n^{-2})$, the required upper tail probability is of order $1/n$, and the Gaussian tail estimate $1-\Phi(\theta)\sim\phi(\theta)/\theta$ gives $\theta\sim\sqrt{2\log n}$.
\end{proof}

\begin{remark}
Independence is a strong assumption. Repeated challenges may reuse evidence, be correlated, or meet preclusion doctrines, and the exact expression $1-\Phi(\theta)^n$ does not survive these complications. The qualitative conclusion does. If challenge evidence is equicorrelated Gaussian with correlation $\varrho\in[0,1)$, write $\epsilon_i=\sqrt{\varrho}\,Z_0+\sqrt{1-\varrho}\,Z_i$ with independent standard normal $Z_0,Z_1,\dots$; then $\max_{i\leq n}\epsilon_i=\sqrt{\varrho}\,Z_0+\sqrt{1-\varrho}\,\max_{i\leq n}Z_i\to\infty$ almost surely, so any fixed threshold is eventually crossed with probability one. Correlation slows attrition without stopping it. Only perfectly correlated challenges, which amount to a single challenge, escape it.
\end{remark}

Resources determine $n$. A well-funded party can litigate, lobby, petition, and reapply repeatedly; a party dependent on the settlement typically cannot. Under a symmetric fixed threshold, settlements that protect weaker parties decay under repeated challenge by stronger ones, not because the evidence ever warranted reversal, but because enough draws from noise eventually cross any fixed line. Symmetric reopening rules therefore function asymmetrically. The formal remedy follows from \eqref{eq:threshold-n}: the threshold for reopening a settlement should rise with the number of prior challenges by the same or allied parties, which is the logic of finality doctrines such as res judicata, or the cost of each challenge should be borne by the challenger in proportion to the reliance interests it threatens. Both remedies need care. Counting prior challenges requires identity rules that treat allied or affiliated challengers as one, so that a party cannot reset the count through nominally separate entities. Evidence of a kind unavailable to earlier challenges should reopen the count. And cost-shifting that is not scaled to the challenger's resources suppresses legitimate challenges by weaker parties as effectively as it deters repeated ones by stronger parties.

The same analysis supports the opposite asymmetry for correction of genuine errors. A party harmed by a mistaken closure may be able to mount only one challenge, and for that party a threshold calibrated to resist a well-resourced repeat challenger is too high. Responsible incompletion therefore requires thresholds that depend on direction and history: low for the first challenge by those who bear a closure's harms, high and rising for repeated challenges against arrangements on which others rely. This is not a compromise between openness and stability. It is the application of the distributive condition of Definition~\ref{def:public} to the procedure of reopening itself.

%======================================================================
\section{Asymmetric Completion}\label{sec:asymmetric}
%======================================================================

The mathematical and normative arguments can now be assembled into a single decision rule. For each component $j$ of an artifact or institution, let $E_j$ be its enabling value (Section~\ref{sec:selective}), $F_j$ the value of the identified alternatives its closure removes, $R_j$ its reversal cost, and $\sigma_j$ the uncertainty surrounding its future use. The expected reversal burden $\bar R_j$ of Section~\ref{sec:selective} is the probability that closure proves mistaken times $R_j$; linearizing that probability in uncertainty gives $\bar R_j\approx\lambda_RR_j\sigma_j$ with $\lambda_R\geq 0$. Because $F_j$ is restricted to alternatives identifiable at closure, while $\lambda_RR_j\sigma_j$ carries the risk of error under unresolved uncertainty, the two terms do not count the same loss twice. Finally let
\[
M_j^+\ =\ \max\{0,\ \Delta\mathcal{M}_j\}
\]
be the increase, if any, in authority--affectedness divergence that closing $j$ under the current authority matrix would produce (Section~\ref{sec:political}). Closures that reduce the divergence receive no penalty on this term, and any such benefit belongs in $E_j$. Define the \emph{asymmetric weight}
\begin{equation}\label{eq:asym-weight}
\tilde w_j\ =\ E_j-F_j-\lambda_RR_j\sigma_j-\lambda_DM_j^+,
\end{equation}
where $\lambda_R\geq 0$ and $\lambda_D\geq 0$ express the normative weight assigned to irreversible error and to unequal continuation rights. With $\lambda_D=0$ the weight reduces to $w_j$ of Section~\ref{sec:selective}. Both $E_j$ and $F_j$ are aggregates over the affected population, of the form $\sum_iw_i(\cdot)$, computed under the controller who would exercise the relevant continuations. The criterion therefore evaluates closure by its consequences for $\sum_iw_iO_i^j$ rather than for the controller alone.

\begin{definition}[Closure criterion]
Component $j$ satisfies the closure criterion if
\begin{equation}\label{eq:closure-criterion}
E_j\ >\ F_j+\lambda_RR_j\sigma_j+\lambda_DM_j^+ .
\end{equation}
\end{definition}

The criterion is exact for purpose vertices and must be generalized for foundations. By Corollary~\ref{cor:purposes}, a sink belongs to an optimal closure only if $\tilde w_j\geq 0$; by Corollary~\ref{cor:foundations}, a foundation may be closed despite failing \eqref{eq:closure-criterion} individually, provided the descendant-closed set it supports satisfies the criterion in aggregate. The general rule is therefore: \emph{close the minimal maximum-weight ideal of the dependency graph under the weights $\tilde w$}. Proposition~\ref{prop:ideal} guarantees that this set exists, is unique, and is computable. Choosing the minimal rather than the maximal optimum breaks ties toward openness. This is justified by an asymmetry of error when, and only when, delayed closure remains feasible: a component wrongly left open can then be closed later at the cost of delay, whereas one wrongly closed can be reopened only at cost $R_j$, if at all. Some opportunities to close also expire. A treaty window, a coordinated standard, a temporary safety regime, or an ecological intervention may be available only for a limited time, and for such components the asymmetry reverses. The tie-breaking rule therefore favors openness only where delayed closure remains available at acceptable cost.

\begin{proposition}[Monotone comparative statics]\label{prop:comparative}
Let $S_\ast(\lambda_R,\lambda_D)$ be the minimal optimal ideal under weights \eqref{eq:asym-weight}. Then $S_\ast$ is nonincreasing in $\lambda_R$ and in $\lambda_D$: raising the weight on irreversible error or on unequal continuation rights never enlarges the set of closed components.
\end{proposition}

\begin{proof}
Write $\tilde w_j(\lambda)=b_j-\lambda c_j$ for $\lambda=\lambda_R$, with $c_j=R_j\sigma_j\geq 0$ and $\lambda_D$ held fixed. Let $\lambda<\lambda'$, $S=S_\ast(\lambda)$, $S'=S_\ast(\lambda')$, and write $J_\lambda$ for the objective. By modularity,
\[
J_{\lambda'}(S')-J_{\lambda'}(S\cap S')=\sum_{j\in S'\setminus S}(b_j-\lambda'c_j)\ \leq\ \sum_{j\in S'\setminus S}(b_j-\lambda c_j)=J_\lambda(S\cup S')-J_\lambda(S)\ \leq\ 0,
\]
the last inequality by optimality of $S$ at $\lambda$. Hence $S\cap S'$, an ideal, is optimal at $\lambda'$, and minimality of $S'$ gives $S'\subseteq S\cap S'\subseteq S$. The argument for $\lambda_D$ is identical with $c_j=M_j^+\geq 0$, which is where the nonnegativity of $M_j^+$ is needed.
\end{proof}

The proof is the same argument as Lemma~\ref{lem:nested}, applied to a normative parameter rather than to time. The proposition shows that the framework's normative parameters act in the expected direction without any further assumption about the dependency structure. Greater concern with irreversibility, or with the concentration of revision authority, contracts closure toward the foundations whose enabling value is great enough to survive the heavier penalties.

\subsection{What must be closed}

The rule is called asymmetric because it treats two kinds of component differently. Some transformations destroy the conditions of participation if left open. A structural system that any occupant could modify would endanger every other occupant; an interface whose semantics any implementer could change would break every other implementation; a safety standard open to case-by-case waiver would protect no one. For such components the enabling value $E_j$ includes the protection of everyone else's participation, and it typically dwarfs the foreclosure and reversal terms. They satisfy the criterion and should be closed firmly, with high reopening thresholds.

Other transformations determine purposes that present authorities cannot legitimately settle in advance: how a space will be used, what a record must be able to say, which ends an assistive system serves, how a community will allocate what it shares. For these, enabling value is small, uncertainty is high, reversal is costly because of the increasing returns discussed in Section~\ref{sec:option}, and closure by a narrow authority raises $\mathcal{M}$. They fail the criterion and should remain provisionally open, or be equipped with inexpensive reopening procedures and continuation rights for those they affect.

\subsection{The declaration requirement}

Exact numerical estimation of the terms in \eqref{eq:asym-weight} is usually impossible, and a framework that demanded it would be useless. The principal function of the criterion is not calculation but \emph{declaration}. Every closure should be accompanied by a record stating its enabling benefit, the alternatives it removes, the burden of reversing it, the uncertainty surrounding its future use, the distribution of authority to revise it, the threshold and procedure for reopening, and the steward responsible. Such a record converts a closure from an event into a revisable claim. It raises $m_L$ in Observation~\ref{obs:measurability}, supplies the misfit registry of Section~\ref{sec:institutions}, keeps $\mathrm{doc}_e$ high in Proposition~\ref{prop:ruin}, and fixes the reopening thresholds whose calibration Proposition~\ref{prop:attrition} showed to be consequential. Even ordinal judgments suffice for much of this. It is often clear that one component's reversal cost exceeds another's, or that one closure affects a population with no authority over it, without any cardinal estimate of either.

%======================================================================
\section{Objections and Boundary Conditions}\label{sec:objections}
%======================================================================

The thesis is best tested against the strongest objections. Each is answered here by narrowing rather than abandoning the argument.

\paragraph{Indecision and diffuse responsibility.} Retained openness, it is objected, produces endless consultation, postponed decisions, and responsibility that belongs to no one. There is also a cognitive version of the objection: unfinished tasks occupy the mind, and a life or institution full of open loops is exhausting. The reply is that the costs described attach to illegible openness, not to openness as such. Masicampo and Baumeister found that the intrusive cognitive effects of unfulfilled goals were largely eliminated when participants formed specific plans for pursuing them, even though the goals themselves remained unfulfilled \cite{masicampo2011}. What relieved the burden was not completion but a legible continuation: a known next step. The open case of Section~\ref{sec:institutions}, with its named steward and stated closing conditions, is the institutional analogue. Latent incompletion as defined in Section~\ref{sec:inert} requires legibility; it is precisely not an unowned loop. Where openness does produce paralysis, the stopping condition \eqref{eq:stopping} shows that deferral has overrun its value and closure is warranted.

\paragraph{Exploitation by hostile agents.} Attachment points can be used to attack as well as to extend. A documented interface is also an attack surface, and an appeal procedure can be used to harass. The reply is that protection against hostile use is a condition of participation and therefore falls among the components the closure criterion directs us to close. Admission control, authentication, rate limits, and the rising reopening thresholds of Proposition~\ref{prop:attrition} are closures whose enabling value consists in protecting everyone else's ability to continue. The thesis does not require that every interface be open to every agent. It requires that the closures be justified by the protection they provide rather than by the proprietor's convenience, and that they not quietly reduce the accessibility coefficient of legitimate continuers to zero.

\paragraph{Coordination at scale.} Large systems require standards, and standards are closures. Without a fixed gauge, railways cannot connect; without fixed protocols, networks cannot interoperate. The reply accepts the premise and locates the closure. Arthur's analysis of increasing returns shows that standardization under competition can lock in designs that are inferior but early \cite{arthur1989}, which strengthens the case for being careful about \emph{what} is standardized. Proposition~\ref{prop:portfolio} and the design-rule framework of Baldwin and Clark \cite{baldwinclark2000} indicate the answer: standardize interfaces, which are foundations with high enabling centrality, and leave implementations and uses open. A standard gauge constrains trains, not destinations.

\paragraph{Misuse by future agents.} The options preserved for successors may be used badly. Future agents may make worse choices than present ones would have made. This is true and is not an objection to preservation. A preserved option is not an endorsed choice; it is a retained capacity to choose, and Proposition~\ref{prop:option} measures its value by the best use available to an aligned chooser, not by the worst. Proposition~\ref{prop:whose-option} states what changes when the chooser is not aligned. Where specific future uses would be catastrophic, they are excluded from the admissible set $\mathfrak{A}$ before breadth or option value is computed, so their preservation never counts in the argument's favor. Where uses are harmful without being inadmissible, the control-indexed valuation handles them: the foreclosure cost of removing an option from a controller misaligned with those affected can be negative, and Proposition~\ref{prop:whose-option} states when. The objection has force only against a thesis of maximal openness, which this essay rejects, not against the claim that present authorities should not foreclose what they lack the knowledge or standing to settle.

\paragraph{Valuing the counterfactual.} It may be objected that foreclosed possibilities cannot be valued without speculative counterfactuals, so that the entire accounting rests on imagined losses. The reply has three parts. First, many foreclosures are not speculative at all: the alternative floor plan, the omitted database field, and the alternative suggestion are concrete and describable at the moment of closure. Second, the framework does not require point estimates. Propositions~\ref{prop:ideal} and~\ref{prop:comparative} operate with any weights, and much of the argument proceeds from inequalities and orderings that are robust to large errors in magnitude. Third, the alternative to valuing counterfactuals imperfectly is to value them at zero, which Observation~\ref{obs:measurability} shows to be a systematic bias rather than a neutral default. An imperfect ledger of possibilities is better than an accounting system in which possibilities cannot appear.

\paragraph{Political feasibility.} Asymmetric completion asks those who hold concentrated revision authority to share it or to carry reopening obligations, and they are the parties best placed to resist. Proposition~\ref{prop:revision-capacity} predicts the form of resistance: added flexibility will be offered, but in rows of the authority matrix that are already full. The framework does not dissolve this problem. It identifies the least demanding point of entry, which is the declaration requirement of Section~\ref{sec:asymmetric}. Requiring that closures be recorded with their foreclosed alternatives and reopening terms redistributes no authority directly, yet it raises the measurability of foreclosure and gives later claimants something to cite.

\paragraph{Open problems.} Several questions are raised by the framework and not settled by it. First, the closure criterion \eqref{eq:closure-criterion} compares terms that ordinal judgment cannot always aggregate. Ordinal information decides a case when it establishes dominance, for instance when a lower bound on $E_j$ exceeds an upper bound on the sum of the penalty terms, or the reverse. Between those bounds lies an indeterminate region in which the criterion does not decide, and there it functions as a requirement that the competing estimates be declared and deliberated rather than as a rule. Second, the declaration requirement leaves institutional questions open: who is responsible for producing a closure record, who has standing to challenge it, and what follows when a record is incomplete or misleading. A plausible default assigns production to the closing authority, standing to anyone whose affectedness exceeds a threshold in $\mathsf{A}$, and treats a materially incomplete record as lowering the reopening threshold of Section~\ref{sec:reopening}, but these are proposals, not results. Third, the admissible set $\mathfrak{A}$ is formally clean and normatively thin. The framework does not say how catastrophic states are identified in advance, who may declare a state catastrophic, or how disagreement about catastrophe is to be resolved, and these are precisely the cases where closure matters most. Treating the declaration of $\mathfrak{A}$ as itself a closure, subject to the same criterion one institutional level up, is the natural extension, but it is not developed here. Fourth, the framework represents departures from proportional authority as reweightings of $Q$ and $\mathsf{A}$ without supplying a theory of which departures are justified. Finally, the domain models share a core structure but not a common dynamics or a common valuation. A unified quantitative theory would require both: a single account of how continuation sets evolve under closure, decay, and maintenance across domains, and a valuation commensurable across the agents and institutions those domains contain. Whether such a theory is possible, or whether continuation is irreducibly plural across domains, is left open.

\paragraph{Boundary conditions.} The thesis therefore does not defend maximal openness, universal reversibility, or indefinite postponement. It defends closure at the minimum level required to secure safety, accountability, interoperability, and dependable participation, combined with retained revisability wherever present authorities lack either the knowledge or the standing to settle future purposes. Where uncertainty is negligible, reversal is cheap, or no meaningful alternatives exist, the diagnostic \eqref{eq:Lambda} is small and ordinary completion is appropriate. Such domains are easy to name: an emergency repair to a failing structure, a crossing people need now, the correction of a known error in a record, the like-for-like replacement of a worn component, and in general any case where the costs of delay are immediate and the foreclosed alternatives are few or well understood. The argument applies with full force only where the three conditions of uncertainty, irreversibility, and preserved diversity coincide. They coincide more often than conventional accounting admits.

%======================================================================
\section{Conclusion: Finishing Without Exhausting the Future}
%======================================================================

This essay began with the observation that an unfinished thing is usually understood through the completed thing it has failed to become. The analysis reverses that relation. A completed thing can also be understood through the continuations it has made impossible. Completion, on this view, is not the disappearance of deficiency but the exercise of an option; not a property of an object but a relation among a state, an authority, a time, and a declared scope; not a single terminal event but a schedule of commitments over a dependency structure.

The formal results support a consistent picture. Completion can raise present fitness while contracting the differentiated future an artifact retains. Accounting that stops at the declaration boundary rewards the displacement of costs onto successors, and accounting that measures achievements more reliably than foreclosures systematically overselects closure. Neither total closure nor total openness maximizes continuability; the optimum is an interior ideal of the dependency graph, firm at foundations whose enabling value repays their cost and open at purposes whose closure no present authority can justify. The same structure governs the plasticity of developing agents, the design of interfaces, the influence of predictive assistance, the representational capacity of institutions, and the allocation of political authority. In several of these domains, premature closure produces the evidence cited in its own defense. In the political domain the analysis also separates two things usually run together: consultation, which improves the information on which an authority closes, and continuation rights, which determine who may act after it has closed. Only the second distributes the capacity to continue.

The framework has a natural companion in the theory of admissible degradation. That account asks which invariants must survive when capacity contracts; the present one asks which alternatives must survive when a project declares success. The first protects continuation against too little remaining structure, the second against too much imposed structure. Together they state two orthogonal requirements, continuation under loss and continuation after success. A system that meets only the first survives shocks but cannot be repurposed; a system that meets only the second can be repurposed but may not survive stress. Degradation and completion are the two ways in which continuation is threatened, and the two accounts are complementary defenses of it.

Completion thus emerges as a distribution of temporal authority. Present agents select which questions become infrastructure and which remain available for later answer. A thing is finished well when it can perform its present work without pretending to exhaust its future. Its boundaries should be firm where failure would destroy the conditions of continuation and permeable where closure would convert participation into obedience. The unfinished artifact is not valuable because its makers lacked resolve, but because they recognized that their last justified decision was not identical to the world's last possible need.
